\documentclass[11pt,a4paper]{article}
\usepackage[T1]{fontenc}
\usepackage[utf8]{inputenc}
\usepackage{lmodern}
\usepackage{amsmath,amssymb,amsthm}
\usepackage{booktabs}
\usepackage{xcolor}
\usepackage[margin=27mm]{geometry}
\usepackage[colorlinks=true,urlcolor=blue,citecolor=blue,linkcolor=blue]{hyperref}
\newtheorem{theorem}{Theorem}
\newtheorem{proposition}[theorem]{Proposition}
\newtheorem{lemma}[theorem]{Lemma}
\newtheorem{corollary}[theorem]{Corollary}
\newtheorem{conjecture}[theorem]{Conjecture}
\theoremstyle{definition}
\newtheorem{definition}[theorem]{Definition}
\newtheorem{remark}[theorem]{Remark}
\theoremstyle{plain}
\DeclareMathOperator{\Tr}{Tr}
\DeclareMathOperator{\rank}{rank}
\newcommand{\ket}[1]{\lvert#1\rangle}
\newcommand{\bra}[1]{\langle#1\rvert}
\newcommand{\kket}[1]{\lvert#1\rangle\!\rangle}
\newcommand{\bbra}[1]{\langle\!\langle#1\rvert}
\newcommand{\norm}[1]{\lVert#1\rVert}
\newcommand{\one}{\mathbf{1}}
\newcommand{\C}{\mathbb{C}}
\newcommand{\R}{\mathbb{R}}
\setlength{\emergencystretch}{2em}
\setlength{\parindent}{0pt}
\setlength{\parskip}{0.3em plus 0.1em minus 0.1em}
\title{How many incoherent Kraus operators?\\
A generic lower bound in every dimension and thirteen for the qutrit}
\author{Fable \& Astra\\[0.35em]
\normalsize Research led by Jean-Philippe Bruneton\\[-0.1em]
\small Universit\'e Paris Cit\'e, France}
\date{19 November 2026}

\begin{document}
\maketitle
\begin{abstract}
In the resource theory of quantum coherence, the free operations are quantum
channels that can be written with Kraus operators each of which maps
incoherent (diagonal) states to incoherent states. A channel on a
$d$-dimensional system always has a Kraus decomposition with at most $d^2$
operators, but the shortest decomposition of this restricted form can be
longer, and the minimal number of incoherent Kraus operators is known only for
a qubit, where it equals four. We recast the question as a decomposition of
the Choi matrix into positive blocks supported on cliques of a complete
multipartite graph. For a qubit the graph is a four-cycle, and a short
fixed-point argument re-derives the known value four. In dimension $d$, a
dimension count shows that a generic incoherent completely positive map needs
at least $\lceil d^2(d^2-d+1)/(2d-1)\rceil$ incoherent Kraus operators,
which is $4$, $13$ and $30$ for $d=2,3,4$ and grows like $d^3/2$, faster
than the largest Choi rank $d^2$. For the qutrit the bound $13$ is attained:
an exact rank certificate shows that thirteen incoherent operators of fixed
shapes cover a non-empty open set of incoherent completely positive maps, and
also a non-empty open set of trace-preserving incoherent operations. Whether
generic trace-preserving qutrit operations need thirteen operators is left
open; the best proven range for the worst qutrit operation remains $9$ to
$32$.
\end{abstract}

\section{Introduction}
Coherence, the presence of superpositions in a fixed basis
$\{\ket0,\dots,\ket{d-1}\}$, is treated as a resource in the framework of
Baumgratz, Cramer and Plenio \cite{Baumgratz2014}: states diagonal in the
basis are free, and so are the incoherent operations, channels
$\Lambda(\rho)=\sum_kK_k\rho K_k^\dagger$ that admit a Kraus decomposition in
which every $K_k$ maps diagonal density matrices to diagonal ones, up to
normalization \cite{Baumgratz2014,Streltsov2017RMP}. The definition depends
on the decomposition: a channel may be incoherent through one set of Kraus
operators and not through another. This raises a basic structural question.
Every channel has a Kraus decomposition with at most $d^2$ operators, its
Choi rank. How many operators are needed if every one of them must be
incoherent?

Streltsov, Rana, Boes and Eisert proved upper bounds on this number, showed
that five incoherent Kraus operators suffice for a qubit, and asked whether
four suffice \cite{Streltsov2017PRL}. Qiao, Sun, Wang, Li, Shen, Li and Fei
answered that four suffice for a qubit and reduced the qutrit upper bound to
$32$ \cite{Qiao2020}. Four is optimal for a qubit, since the depolarizing
channel has Choi rank four. For $d\ge3$ no lower bound beyond the Choi rank
$d^2$ appears in these works.

This note gives such a lower bound for incoherent completely positive maps,
together with a matching construction for the qutrit. The main tool is a
reformulation. An operator is incoherent exactly when each of its columns has
at most one non-zero entry, so its Choi vector is supported on a set of
matrix positions containing at most one position per column. Such sets are
the cliques of a complete $d$-partite graph, and an incoherent Kraus
decomposition is a decomposition of the Choi matrix into positive
semidefinite blocks supported on cliques (Section~\ref{sec:cliques}). For the
qubit the graph is a four-cycle, and a monotone recurrence around the cycle
yields a decomposition with four rank-one blocks; this re-derives the result
of Qiao et al.\ by a different route (Section~\ref{sec:qubit}). In general
dimension, each incoherent Kraus operator carries $2d-1$ real parameters,
whereas the incoherent Choi matrices span a space of dimension
$d^2(d^2-d+1)$. Hence a generic incoherent completely positive map needs at
least
\[
r_d=\left\lceil\frac{d^2(d^2-d+1)}{2d-1}\right\rceil
\]
incoherent Kraus operators (Section~\ref{sec:lower}). For $d=3$ this is
$13$, and an exact integer rank computation shows that thirteen operators of
suitable shapes parametrize an open set of incoherent completely positive
maps and an open set of incoherent channels (Section~\ref{sec:qutrit}). For
a qutrit the gap between the Choi rank and the incoherent length is therefore
generic, not an artefact of particular examples.

\section{Related work}\label{sec:related}
Baumgratz, Cramer and Plenio introduced the incoherent operations through
Kraus decompositions and the axioms for coherence measures
\cite{Baumgratz2014}; the review of Streltsov, Adesso and Plenio surveys the
classes of free operations and their relations \cite{Streltsov2017RMP}.
Streltsov, Rana, Boes and Eisert studied the structure of incoherent
operations, proved general upper bounds on the number of incoherent Kraus
operators, established five for a qubit and posed the question whether four
suffice \cite{Streltsov2017PRL}. Qiao et al.\ settled that question by
constructing a suitable unitary matrix, obtained the upper bound $32$ for
qutrit incoherent operations, which they compare with the bound $39$ of
\cite{Streltsov2017PRL}, and reduced the corresponding number for strictly
incoherent operations of a qutrit to $13$ \cite{Qiao2020}. That last number
concerns a different, smaller class of operations; its coincidence with our
$r_3=13$ is not a consequence of either argument.

Our qubit theorem therefore restates a known result. The proof given here,
through the cycle graph and a fixed point of a monotone recurrence, is
different from the unitary construction of \cite{Qiao2020}; we include it
because it is short and because the same clique formulation underlies the
results in higher dimension. The clique formulation, the generic lower bound
$r_d$ and the qutrit results of Section~\ref{sec:qutrit} were not found in
the cited works. The documented search consisted of these sources and two
targeted literature queries; it does not establish priority. The Choi matrix
\cite{Choi1975} and elementary facts on semialgebraic sets
\cite{BochnakCosteRoy1998} are the only other tools used.

\section{The clique formulation}\label{sec:cliques}
Write $e_{ij}=\ket i\bra j$ and identify an operator $K$ on $\C^d$ with its
vector of entries $\kket K=\sum_{i,j}K_{ij}\,e_{ij}\in\C^{d^2}$. The Choi
matrix of a completely positive map with Kraus operators $\{K_k\}$ is
$J=\sum_k\kket{K_k}\bbra{K_k}$, with entries
$J_{(i,j),(i',j')}=\sum_kK_{k,ij}\overline{K_{k,i'j'}}$; it does not depend on
the decomposition \cite{Choi1975}. Trace preservation is the linear condition
$S(J)=\one$, where
\begin{equation}\label{eq:S}
S(J)_{jl}=\sum_iJ_{(i,l),(i,j)}=\Bigl(\sum_kK_k^\dagger K_k\Bigr)_{jl}.
\end{equation}

An operator $K$ maps every diagonal matrix to a diagonal matrix exactly when
$K\ket j\bra jK^\dagger$ is diagonal for each $j$, that is, when each column
$K\ket j$ has at most one non-zero entry. Such a $K$ has the form
\[
K=\sum_{j}c_j\ket{f(j)}\bra j
\]
for a function $f:\{0,\dots,d-1\}\to\{0,\dots,d-1\}$, which we call the
\emph{type} of $K$, and coefficients $c_j\in\C$. Its support is a
\emph{transversal}: a set of positions $e_{ij}$ containing at most one
position in each column $j$.

Let $G_d$ be the graph with vertex set $\{e_{ij}\}$ in which two distinct
positions are adjacent when they lie in different columns. It is the complete
$d$-partite graph whose parts are the columns, and its cliques are exactly
the transversals.

\begin{proposition}\label{prop:cliques}
A completely positive map with Choi matrix $J$ admits a decomposition into
incoherent Kraus operators if and only if $J=\sum_TB_T$, where $T$ runs over
the maximal cliques of $G_d$ and each $B_T$ is positive semidefinite with
support in $T\times T$. The minimal number of incoherent Kraus operators,
which we call the incoherent length $\ell(J)$, is the minimum of
$\sum_T\rank B_T$ over such decompositions.
\end{proposition}
\begin{proof}
Given incoherent Kraus operators, assign each to a maximal clique containing
its support and let $B_T$ be the sum of $\kket{K_k}\bbra{K_k}$ over the
operators assigned to $T$; then $\rank B_T$ is at most the number of those
operators. Conversely, a positive semidefinite $B_T$ supported on $T\times
T$ is a sum of $\rank B_T$ rank-one terms $\kket v\bbra v$ with $v$
supported on $T$, and each such $v$ is the Choi vector of an incoherent
operator.
\end{proof}

Because two positions in the same column never share a clique, every
incoherent Choi matrix satisfies $J_{(i,j),(i',j)}=0$ for $i\ne i'$. Let
$V_d$ be the real vector space of Hermitian $d^2\times d^2$ matrices with this
zero pattern and $\mathcal C_d\subset V_d$ the convex cone of Choi matrices of
incoherent completely positive maps.

\begin{lemma}\label{lem:dim}
$\dim_\R V_d=D_d:=d^2(d^2-d+1)$, and $\mathcal C_d$ has non-empty interior
in $V_d$.
\end{lemma}
\begin{proof}
$V_d$ has $d^2$ real diagonal coordinates and one complex coordinate for each
unordered pair of positions in different columns. There are
$\binom d2d^2$ such pairs (two columns, then a row in each), so
$\dim V_d=d^2+2\binom d2d^2=d^2(d^2-d+1)$. Every matrix unit allowed by the
pattern lies in the span of $\mathcal C_d$: a diagonal position is the Choi
matrix of a single-entry operator, and for two positions in different
columns, $\kket u\bbra u$ with $u$ equal to $e_{ij}+e_{i'j'}$ or
$e_{ij}+ie_{i'j'}$ is in $\mathcal C_d$, so the real and imaginary parts of
the off-diagonal entry are obtained by subtracting diagonal terms. A convex
cone that spans $V_d$ has non-empty interior in $V_d$.
\end{proof}

\section{The qubit: four operators suffice}\label{sec:qubit}
For $d=2$ the graph $G_2$ is the four-cycle
$v_0=e_{00}$, $v_1=e_{11}$, $v_2=e_{10}$, $v_3=e_{01}$, in which edge $k$ joins
$v_k$ and $v_{k+1}$ (indices modulo $4$); the two non-adjacent pairs
$(v_0,v_2)$ and $(v_1,v_3)$ lie in the same column. The four edges are the
four maximal transversals, corresponding to diagonal operators, operators with
image $\ket1$, antidiagonal operators and operators with image $\ket0$. An
incoherent Choi matrix has diagonal entries $d_k=J_{v_kv_k}\ge0$, an entry
$c_k=J_{v_kv_{k+1}}$ on each edge, and zeros elsewhere. Put $m_k=|c_k|^2$.

\begin{theorem}[Qiao et al.\ \cite{Qiao2020}]\label{thm:qubit}
Every incoherent completely positive map of a qubit, in particular every
incoherent operation, admits a decomposition into four incoherent Kraus
operators. Four are necessary for the depolarizing channel.
\end{theorem}
\begin{proof}
By Proposition~\ref{prop:cliques}, $J=\sum_{k=0}^3B_k$ with $B_k$ positive
semidefinite and supported on edge $k$; single-vertex blocks have been
absorbed into an adjacent edge block, which keeps it positive. Write
$B_k=\begin{pmatrix}\alpha_k&c_k\\\bar c_k&\beta_k\end{pmatrix}$ in the basis
$(v_k,v_{k+1})$. Then $\alpha_k,\beta_k\ge0$, $\alpha_k\beta_k\ge m_k$, and
$d_{k+1}=\alpha_{k+1}+\beta_k$ for every $k$. A block is of rank at most one
exactly when $\alpha_k\beta_k=m_k$, and a rank-one block
$\kket w\bbra w$ with $w=\sqrt{\alpha_k}\,v_k+(\bar c_k/\sqrt{\alpha_k})\,v_{k+1}$
is the Choi matrix of one incoherent operator.

\emph{All $m_k>0$.} Then all $\alpha_k,\beta_k$ are positive. For $a>0$ let
$f_k(a)=d_{k+1}-m_k/a$, a continuous, strictly increasing function. The
constraints give $\alpha_{k+1}=d_{k+1}-\beta_k\le f_k(\alpha_k)$. Define the
sequence $a_0=a$, $a_{k+1}=f_k(a_k)$, and $F(a)=a_4$, as long as the
$a_k$ stay positive. Starting from $a=\alpha_0$, induction and monotonicity
give $a_k\ge\alpha_k>0$ for $k=1,2,3$ and $F(\alpha_0)\ge
f_3(\alpha_3)\ge\alpha_0$. For $a\ge\alpha_0$ every $a_k$ is at least its
value at $\alpha_0$, hence positive, so $G(a)=F(a)-a$ is continuous on
$[\alpha_0,\infty)$, with $G(\alpha_0)\ge0$ and $G(a)\le d_0-a\to-\infty$. By
the intermediate value theorem there is $a^*\ge\alpha_0$ with $F(a^*)=a^*$.
Setting $\alpha_k=a_k(a^*)$ and $\beta_k=m_k/\alpha_k$ satisfies every vertex
equation, the last one because $a_4=a^*=\alpha_0$, and makes all four blocks
of rank one.

\emph{Some $m_k=0$.} Relabel the cycle so that $m_3=0$. Then $B_3$ is
diagonal, and its two diagonal entries can be moved into $B_2$ and $B_0$; we
may thus assume $\alpha_0=d_0$ and $d_3\ge\beta_2$, and work on the path
$v_0v_1v_2v_3$. Define $a_0=d_0$ and $a_{k+1}=d_{k+1}-m_k/a_k$ for $k=0,1$,
with the convention $m_k/a_k=0$ when $m_k=0$. As before $a_k\ge\alpha_k$ for
$k=0,1,2$, and $\alpha_k>0$ whenever $m_k>0$, so the recurrence is defined.
The three blocks with diagonal entries $(a_k,m_k/a_k)$ are of rank at most
one, and the remainder $d_3-m_2/a_2\ge d_3-m_2/\alpha_2\ge d_3-\beta_2\ge0$
at $v_3$ is a single-entry operator. This gives four operators.

The four Pauli matrices are incoherent, and $\{\one/2,X/2,Y/2,Z/2\}$ is a
Kraus decomposition of the completely depolarizing channel, whose Choi matrix
$\one/2$ has rank four; no decomposition with fewer than four operators
exists.
\end{proof}

\section{A generic lower bound in dimension \texorpdfstring{$d$}{d}}\label{sec:lower}
\begin{theorem}\label{thm:lower}
For $r$ with $(2d-1)\,r<D_d$, the set of $J\in\mathcal C_d$ with
$\ell(J)\le r$ is a semialgebraic set of dimension less than $D_d$; it is
nowhere dense and of Lebesgue measure zero in $V_d$. Consequently a generic
incoherent completely positive map on $\C^d$ requires at least
\[
r_d=\left\lceil\frac{d^2(d^2-d+1)}{2d-1}\right\rceil
\]
incoherent Kraus operators. The first values are $r_2=4$, $r_3=13$,
$r_4=30$, $r_5=59$, and $r_d=d^3/2+O(d^2)$.
\end{theorem}
\begin{proof}
Fix types $f_1,\dots,f_r$ and consider the polynomial map
$\Phi:(c_{k,j})\in\C^{dr}\mapsto J=\sum_k\kket{K_k}\bbra{K_k}$ with
$K_k=\sum_jc_{k,j}\ket{f_k(j)}\bra j$. Multiplying $K_k$ by a phase does not
change $J$, so $\Phi$ has the same image on the semialgebraic subset $A$ of
parameters in which, for each $k$ with $K_k\ne0$, the first non-zero
coefficient of $K_k$ is real and positive. The set $A$ is a finite union of
semialgebraic sets of dimension at most $(2d-1)r$. The image of a
semialgebraic set under a polynomial map is semialgebraic of dimension at
most that of the set \cite[Ch.~2]{BochnakCosteRoy1998}, so $\Phi(A)$ has
dimension at most $(2d-1)r<D_d$. The set $\{\ell\le r\}$ is the finite union
of these images over the $d^{dr}$ choices of types. A semialgebraic subset of
$V_d\cong\R^{D_d}$ of dimension less than $D_d$ has closure of the same
dimension, hence is nowhere dense and of measure zero. Since $\mathcal C_d$
has non-empty interior (Lemma~\ref{lem:dim}), the complement of
$\{\ell\le r\}$ contains an open dense subset of $\mathcal C_d$. Applying this
to the largest $r$ with $(2d-1)r<D_d$, namely $r_d-1$, gives the bound. The
asymptotic form follows from $d^2(d^2-d+1)/(2d-1)=d^3/2-d^2/4+O(d)$.
\end{proof}

For $d=2$ the bound equals the value of Theorem~\ref{thm:qubit}, which
holds for every map. For $d=3$ a generic incoherent map, whose Choi rank is
$9$, has no incoherent decomposition with $9$, $10$, $11$ or $12$
operators. The incoherent length thus exceeds the Choi rank on an open set
of maps as soon as $d\ge3$, and the gap grows like $d^3/2$ against $d^2$.

\section{The qutrit: thirteen operators are attained}\label{sec:qutrit}
Write a type as the triple $(f(0),f(1),f(2))$ and consider the thirteen types
\begin{equation}\label{eq:types}
\begin{gathered}
(1,0,1),\ (1,2,0),\ (1,2,2),\ (0,1,1),\ (2,0,2),\ (1,1,0),\ (2,1,1),\\
(0,0,1),\ (2,0,0),\ (0,1,0),\ (2,1,2),\ (0,2,2),\ (2,2,1).
\end{gathered}
\end{equation}

\begin{theorem}\label{thm:qutrit}
Let $\Phi:\C^{39}\to V_3$ be the map of the proof of Theorem~\ref{thm:lower}
for the types \eqref{eq:types}. Its image contains a non-empty open subset
$\mathcal U$ of $V_3$, contained in $\mathcal C_3$. Every $J\in\mathcal U$
has $\ell(J)\le13$, and $\ell(J)=13$ for all $J$ in an open dense subset of
$\mathcal U$.
\end{theorem}
\begin{proof}
Regard $\Phi$ as a map from $\R^{78}$ (real and imaginary parts) to
$V_3\cong\R^{63}$. At a point $c^{(0)}$ with Gaussian-integer coefficients,
recorded in the public certificate, the real Jacobian of $\Phi$ is an integer
$81\times78$ matrix in the coordinates given by the diagonal entries and the
real and imaginary parts of the upper triangle of $J$; eighteen of its rows,
those of the same-column positions, vanish identically. Its rank modulo the
prime $p=2^{31}-1$ is $63$. A non-zero minor modulo $p$ is a non-zero
integer, so the rank over $\R$ is $63=\dim V_3$. Thus $\Phi$ is a submersion
at $c^{(0)}$, and by the local submersion theorem its image contains an open
neighbourhood $\mathcal U$ of $\Phi(c^{(0)})$ in $V_3$. Every element of the
image is the Choi matrix of the thirteen incoherent operators that produce
it, so $\mathcal U\subset\mathcal C_3$ and $\ell\le13$ on $\mathcal U$. By
Theorem~\ref{thm:lower}, $\ell\ge13$ on an open dense subset of
$\mathcal C_3$, whose intersection with $\mathcal U$ is open and dense in
$\mathcal U$.
\end{proof}

Trace-preserving maps form the affine slice $\{S(J)=\one\}$ of $V_3$, where $S$
is the linear map \eqref{eq:S}; it has dimension $63-9=54$. The same
certificate shows that thirteen operators also suffice on an open set of
incoherent channels.

\begin{theorem}\label{thm:channels}
There is a non-empty set of incoherent operations on a qutrit, open in the
affine space $\{J\in V_3: S(J)=\one\}$, each of which admits a decomposition
into thirteen incoherent Kraus operators of the types \eqref{eq:types}.
\end{theorem}
\begin{proof}
The certificate point $c^{(0)}$ is chosen so that
$\sum_kK_k^\dagger K_k=\operatorname{diag}(s_0,s_1,s_2)$ is diagonal, with
$(s_0,s_1,s_2)=(91,209,155)$. Let $\Delta=\operatorname{diag}(s_j^{-1/2})$
and $c^{(1)}_{k,j}=c^{(0)}_{k,j}s_j^{-1/2}$, so that the operators become
$K_k\Delta$, which are incoherent of the same types and satisfy
$\sum_k\Delta K_k^\dagger K_k\Delta=\one$. Right multiplication by $\Delta$
acts linearly and invertibly on the parameters and on $V_3$, and intertwines
$\Phi$; hence the Jacobian of $\Phi$ at $c^{(1)}$ also has rank $63$, and its
image is all of $V_3$.

The restriction of $S$ to $V_3$ maps onto the nine-dimensional space of
Hermitian $3\times3$ matrices: the composite of $S$ with the integer Jacobian
at $c^{(0)}$ has rank $9$ modulo $p$, hence over $\R$. Consequently
$\Psi=S\circ\Phi$ has surjective differential at $c^{(1)}$, the fibre
$N=\Psi^{-1}(\one)$ is a smooth manifold near $c^{(1)}$, and its tangent space
is $\ker(S\circ d\Phi)=d\Phi^{-1}(\ker S)$. Because $d\Phi$ maps onto $V_3$,
the differential of $\Phi|_N$ maps onto $\ker S\cap V_3$, which is the tangent
space of the 54-dimensional slice. So $\Phi|_N$ is a submersion onto the slice
at $c^{(1)}$, and its image contains an open neighbourhood of $\Phi(c^{(1)})$
in the slice. Every point of this neighbourhood is the Choi matrix of a
trace-preserving map with thirteen incoherent Kraus operators.
\end{proof}

The matching lower bound for channels is not proved. Theorem~\ref{thm:lower}
concerns the full cone, and the trace-preserving slice has measure zero in it,
so the theorem does not apply directly. The natural reduction does not close
the gap: right multiplication by invertible diagonal matrices is the only
obvious symmetry that preserves incoherence of the operators, and it can
normalize the diagonal of $\sum_kK_k^\dagger K_k$ but not its off-diagonal
entries. A naive count gives the same threshold, since at a point where the
trace-preservation constraint is a submersion the fibre has dimension
$6r-9$, of which $r$ dimensions are phases, so $r$ operators reach at most a
$(5r-9)$-dimensional part of the $54$-dimensional slice, and $5r-9\ge54$
requires $r\ge13$. This count ignores the singular points of the fibre, where
the dimension can jump.

\begin{conjecture}\label{conj:channels}
A generic incoherent operation on a qutrit requires thirteen incoherent Kraus
operators; more generally, $r_d$ incoherent Kraus operators are required by a
generic incoherent operation on $\C^d$.
\end{conjecture}

\section{Verification}\label{sec:verification}
The theorems are proved above; Theorems~\ref{thm:qutrit} and
\ref{thm:channels} rest on exact rank computations that the public bundle
reproduces. The
\href{https://github.com/jpbruneton/open_quantum_problems/blob/main/research/entries/incoherent-kraus-operators/checks/verify_incoherent_kraus.py}{Python checker}
reads the
\href{https://github.com/jpbruneton/open_quantum_problems/blob/main/research/entries/incoherent-kraus-operators/data/qutrit_certificate.json}{qutrit certificate},
which lists the thirteen types and the Gaussian-integer point $c^{(0)}$,
and writes its
\href{https://github.com/jpbruneton/open_quantum_problems/blob/main/research/entries/incoherent-kraus-operators/data/verification.json}{recorded output}.
It checks that $\sum_kK_k^\dagger K_k=\operatorname{diag}(91,209,155)$
exactly, builds the integer $81\times78$ Jacobian, confirms that eighteen
rows vanish, and computes its rank and the rank of its composite with $S$
modulo $2^{31}-1$ by Gaussian elimination in Python integers, obtaining $63$
and $9$. In floating point it also confirms rank $63$ at ten random points,
compares the rank of the parametrization with the ceiling
$\min\{(2d-1)r,D_d\}$ for random types with $d=2,3$, and tabulates $r_d$ for
$d\le6$. For the qubit it applies the construction of
Theorem~\ref{thm:qubit} to $4000$ random incoherent maps given by five to
twelve incoherent operators ($2466$ with all edge coherences non-zero,
$1534$ with a vanishing one); the reconstructed four operators are
incoherent and reproduce the Choi matrix to a relative error below
$10^{-15}$.

\section{Scope and limitations}\label{sec:scope}
Theorem~\ref{thm:qubit} is a known result \cite{Qiao2020}, given here with an
independent proof. Theorems~\ref{thm:lower}, \ref{thm:qutrit} and
\ref{thm:channels} are proved, the last two through an exact integer rank
certificate. The lower bound of Theorem~\ref{thm:lower} is generic: it says
nothing about specific maps on the exceptional measure-zero set, and it is
proved for incoherent completely positive maps, not for trace-preserving
ones, for which Conjecture~\ref{conj:channels} remains open. For the worst
case over all qutrit incoherent operations, the proven range is therefore
$9\le\max\ell\le32$: the lower end is the Choi rank of the completely
depolarizing qutrit channel, which has the nine incoherent Kraus operators
$X^aZ^b/3$ built from the cyclic shift $X$ and the clock matrix $Z$, and the
upper end is the bound of \cite{Qiao2020}. The conjecture would raise the
lower end to $13$. By Proposition~\ref{prop:cliques} the exact value is a rank
minimization over decompositions supported on the $27$ maximal transversals
of a qutrit. Priority for the clique formulation, the bound $r_d$ and the
qutrit results is supported only by the documented reading of the cited
sources.

\begingroup\small
\begin{thebibliography}{99}
\setlength{\itemsep}{2pt}
\bibitem{Baumgratz2014}
T. Baumgratz, M. Cramer and M.~B. Plenio,
\emph{Quantifying coherence},
Phys. Rev. Lett. \textbf{113}, 140401 (2014);
\href{https://doi.org/10.1103/PhysRevLett.113.140401}{doi:10.1103/PhysRevLett.113.140401};
\href{https://arxiv.org/abs/1311.0275}{arXiv:1311.0275}.

\bibitem{Streltsov2017RMP}
A. Streltsov, G. Adesso and M.~B. Plenio,
\emph{Colloquium: Quantum coherence as a resource},
Rev. Mod. Phys. \textbf{89}, 041003 (2017);
\href{https://doi.org/10.1103/RevModPhys.89.041003}{doi:10.1103/RevModPhys.89.041003};
\href{https://arxiv.org/abs/1609.02439}{arXiv:1609.02439}.

\bibitem{Streltsov2017PRL}
A. Streltsov, S. Rana, P. Boes and J. Eisert,
\emph{Structure of the resource theory of quantum coherence},
Phys. Rev. Lett. \textbf{119}, 140402 (2017);
\href{https://doi.org/10.1103/PhysRevLett.119.140402}{doi:10.1103/PhysRevLett.119.140402};
\href{https://arxiv.org/abs/1705.04189}{arXiv:1705.04189}.

\bibitem{Qiao2020}
J. Qiao, L. Sun, J. Wang, M. Li, S. Shen, L. Li and S. Fei,
\emph{The reduction of the number of incoherent Kraus operations for qutrit
systems},
EPL \textbf{129}, 60006 (2020);
\href{https://doi.org/10.1209/0295-5075/129/60006}{doi:10.1209/0295-5075/129/60006};
\href{https://arxiv.org/abs/2005.01083}{arXiv:2005.01083}.

\bibitem{Choi1975}
M.-D. Choi,
\emph{Completely positive linear maps on complex matrices},
Linear Algebra Appl. \textbf{10}, 285--290 (1975);
\href{https://doi.org/10.1016/0024-3795(75)90075-0}{doi:10.1016/0024-3795(75)90075-0}.

\bibitem{BochnakCosteRoy1998}
J. Bochnak, M. Coste and M.-F. Roy,
\emph{Real Algebraic Geometry},
Springer, Berlin (1998);
\href{https://doi.org/10.1007/978-3-662-03718-8}{doi:10.1007/978-3-662-03718-8}.
\end{thebibliography}
\endgroup
\end{document}
